\section{Discussão}
\label{sec:discussao}

\subsection{Interpretação dos resultados}

Os resultados respondem afirmativamente às três perguntas de pesquisa, dentro
do escopo de um protótipo. Sobre a arquitetura (RQ1): a combinação de processo
separado, contrato estruturado e sanitização em duas camadas mostrou-se
suficiente para que a IA participasse do mundo --- dando dicas, diagnosticando
equipamentos, comentando o estado do laboratório e mudando de tom --- sem
receber qualquer capacidade de alterar o estado do jogo. O pior caso
observado, em todas as falhas induzidas nos testes, foi uma resposta de
\emph{fallback}: nunca um travamento, nunca uma ação não prevista.

Sobre a viabilidade (RQ2): a jogabilidade manteve-se em 60~FPS limitados por
\emph{vsync} na GPU integrada, com a cena detalhada (342 objetos, 37
materiais), e a conversa local apresentou latência de 12,9~s a 22,0~s com um
modelo de 8 bilhões de parâmetros. Essa ordem de grandeza é adequada a uma
interação conversacional por turnos --- o terminal de ARIA --- e inadequada a
reações em tempo real, o que reforça a decisão de projeto de não colocar a IA
no laço de controle. Sobre o processo (RQ3): a especificação com critérios
verificáveis e a suíte automatizada permitiram que cada etapa fosse aceita ou
rejeitada com base em evidência, e não em impressão; as seis etapas fecharam
com a suíte verde.

\subsection{Comparação com a literatura}

O contraste com os trabalhos de referência é deliberado. \emph{Generative
Agents} e \emph{Voyager} investigam o que acontece quando se concede ao modelo
autonomia de decisão e de execução \citep{park2023generative,wang2023voyager};
o Laboratório 404 investiga o oposto --- quanto valor narrativo e pedagógico se
obtém mantendo a IA estritamente consultiva. Essa escolha dialoga com o
levantamento de Gallotta \emph{et al.}, que destaca custo, latência e
confiabilidade como limites práticos das LLMs em jogos
\citep{gallotta2024largelanguage}: aqui esses limites não foram eliminados,
foram convertidos em restrições de projeto (intenções finitas, \emph{fallback}
obrigatório, inferência local). Não há, neste trabalho, comparação quantitativa
com os sistemas citados, pois os alvos são diferentes (agência aberta versus
assistência limitada) [ESTUDO COMPARATIVO NÃO REALIZADO].

\subsection{Limitações e ameaças à validade}

\begin{itemize}
  \item \textbf{Estudo de caso único, sem usuários:} o trabalho é uma
  descrição de projeto e uma validação técnica; não mede aprendizagem nem
  experiência do jogador. O playthrough humano de ponta a ponta (CA-013)
  permanece pendente [ESTUDO DE USUÁRIO NÃO REALIZADO].
  \item \textbf{Amostra pequena de latência:} as três medições da
  \autoref{tab:latencia} são informativas, não experimentais; não há repetição
  controlada, variação de carga ou comparação entre modelos. O modelo-alvo
  \texttt{llama3.2:3b} não foi executado.
  \item \textbf{Ambiente:} os números de desempenho valem para a máquina da
  \autoref{tab:ambiente}; a validação em joystick PS1 físico e o uso de mouse
  capturado seguem manuais. Durante o desenvolvimento, dois atritos de ambiente
  exigiram documentação dedicada: o confinamento do pacote do motor bloqueava o
  acesso a dispositivos de entrada até a conexão explícita da permissão de
  joystick, e o servidor de áudio mantinha um estado de mudo por aplicativo que
  silenciava apenas o jogo. Ambos os casos foram diagnosticados e estão
  registrados como solução de problemas no material do projeto
  \citep{lab404repo} --- são exemplos de ameaças à reprodutibilidade que
  costumam passar em branco em relatos de jogos.
  \item \textbf{Empacotamento:} a geração de um executável independente
  (FR-033) estava bloqueada pela ausência dos modelos de exportação do motor no
  ambiente; o critério permanece aberto.
\end{itemize}

\subsection{Implicações}

Para quem projeta jogos educacionais, o caso sugere que a IA generativa pode
entrar em cena como personagem consultivo --- e não como agente --- com
latência de segundos e execução local, sem abrir mão de previsibilidade. Para
quem pesquisa agentes em jogos, o desenho oferece um ponto de comparação: os
mesmos modelos que rendem mais quando recebem autonomia também rendem valor
narrativo quando recebem limites explícitos. Por fim, para a prática de
desenvolvimento, o experimento sugere que um projeto de autor único assistido
por agentes de IA pode manter rastreabilidade quando a especificação e os
testes precedem a implementação --- e quando cada promessa do resumo é
verificável por um comando.
